\subsection{Переменные}
Переменная — это имя, за которым «закреплено» значение. В Lua есть базовые типы: \hyperlink{term:number}{число}, \hyperlink{term:string}{строка}, \hyperlink{term:boolean}{булево}, \hyperlink{term:table}{таблица}, \hyperlink{term:function}{функция}, \hyperlink{term:nil}{nil}. Узнать тип можно через \hyperlink{term:type}{type()}. \hyperlink{term:scope}{Локальная} переменная объявляется через \verb|local| и видна внутри текущего блока; \textbf{глобальная} — без \verb|local| (видна во всём скрипте).

Число — это величина для арифметики (например, \verb|3.14| или \verb|10|). Строка — текст в кавычках (например, \verb|"Pioneer"|), над ней выполняют конкатенацию \verb|..| и получают длину \verb|#s|. Булево — \verb|true|/\verb|false|, используется в \hyperlink{term:condition}{условиях}. Таблица — универсальная структура: массив \verb|t[1]| и словарь \verb|t.x| одновременно. Функция — именуемый блок кода, например \verb|local function f(x) return x*2 end|.

\begin{minted}{lua}
-- локальная числовая переменная
local x = 10
-- глобальная строковая переменная
name = "Pioneer"
-- локальная булева переменная
local ready = true
-- таблица с позиционными элементами и именованным полем
local t = {1, 2, 3, speed = 3.5}
-- множественное присваивание
x, y = 1, 2
-- обмен значениями
x, y = y, x
-- выясняем типы значений
print(type(x), type(name), type(ready), type(t))  -- number string boolean table
\end{minted}
\begin{minted}{lua}
-- словарь (именованные поля)
local p = {x = 1, y = 2}
-- массив (позиционные элементы)
local a = {10, 20, 30}
-- добавление нового поля
p.z = 3
-- добавление нового элемента массива
a[4] = 40
\end{minted}
\paragraph{Пояснения.}
- Таблица — единственная сложная структура в Lua; через неё строят списки, словари и объекты.
- \verb|nil| при присваивании полю таблицы удаляет это поле: \verb|t.speed = nil|.
- Глобальные переменные удобно использовать только для настроек и состояний, иначе код становится хрупким; предпочитайте \verb|local|.

\subsection*{Задачи для самостоятельного решения}
\begin{enumerate}
  \item Объявите локальную переменную \verb|x| со значением \verb|42|.
  \item Присвойте строку \verb|"Pioneer"| глобальной переменной \verb|name|.
  \item Поменяйте местами значения \verb|a| и \verb|b| одной строкой.
  \item Создайте таблицу \verb|t| с полем \verb|speed=3.5| и элементами \verb|1,2,3|.
  \item Добавьте в таблицу \verb|t| поле \verb|mode| со значением \verb|"auto"|.
  \item Объявите \verb|local flag = false| и затем присвойте глобальной \verb|flag| значение \verb|true|.
  \item Создайте пустую таблицу \verb|buf| и добавьте в неё числа \verb|5| и \verb|7|.
  \item Удалите поле \verb|speed| из таблицы \verb|t|, присвоив \verb|nil|.
  \item Присвойте одновременно значения \verb|r=1|, \verb|g=0|, \verb|b=0|.
  \item Создайте таблицу \verb|color| вида \verb|{r=1,g=0,b=0}|.
  \item Извлеките третий элемент массива \verb|a| в переменную \verb|third|.
  \item Присвойте \verb|local s = "Lua"| и \verb|local n = 3.14|.
  \item Создайте локальную функцию \verb|id|, возвращающую свой аргумент.
  \item Объявите \verb|local config = nil| и затем присвойте \verb|config = {enabled=true}|.
  \item Создайте таблицу \verb|point| и добавьте в неё \verb|x=10| и \verb|y=20|.
\end{enumerate}
